\documentclass[11pt]{article}
\usepackage[margin=1in]{geometry}
\usepackage{amsmath,amssymb,amsthm}
\usepackage{graphicx}
\usepackage{booktabs}
\usepackage{enumitem}

\newtheorem{hypothesis}{Hypothesis}

\title{Supplementary Notes for the Masters Research Report:\\
Figure Captions, Potential Conditions, and Literature Gap Analysis}
\author{Ethan Knox}
\date{March 2026}

\begin{document}
\maketitle


%% ============================================================
\section{Restrictions on the Potential $V$}
%% ============================================================

The following hypotheses on $V:\mathbb{R}^n\to\mathbb{R}$ are drawn from the two primary references.
We organize them by the theorem that requires them.

\subsection{Core Hypotheses (Kirr--Natarajan, 2018, \S1)}

\begin{enumerate}[label=\textbf{(H\arabic*)}]
  \item $V\in L^\infty(\mathbb{R}^n)$.
  \item $\displaystyle\lim_{|x|\to\infty} V(x) = 0$.
  \item $L_0 = -\Delta + V$ has a lowest eigenvalue $-E_0 < 0$.
\end{enumerate}

\noindent\textbf{Remark.}
(H1)--(H2) ensure $V$ is a relatively compact perturbation of $-\Delta$, so that $-\Delta+V$ is self-adjoint on $L^2(\mathbb{R}^n)$ with domain $H^2(\mathbb{R}^n)$ and essential spectrum $[0,\infty)$ (Weyl's theorem). The authors note that (H1) can be relaxed to $V\in L^q(\mathbb{R}^n)$, $q\geq 1$, $q > n/2$, at the cost of complicating certain estimates.

\subsection{For Norm Scaling and Multi-Peak Decomposition (Theorem~4.1)}

In addition to (H1)--(H2), the rescaling analysis requires:
\begin{equation}
  \boxed{x\cdot\nabla V \in L^\infty(\mathbb{R}^n).} \tag{H4$'$}
\end{equation}
This condition controls the Pohozaev-type identity used to derive the scaling laws $\|\psi_E\|_{L^{2p+2}}^{2p+2} \sim E^{1/2+1/p}$ as $E\to\infty$.

\subsection{For Compactness at Finite $E$ (Theorem~3.3)}

The concentration-compactness argument that prevents splitting (profiles drifting to $|x|=\infty$) requires $V\in W^{1,\infty}(\mathbb{R}^n)$ together with the following decay/dominance conditions on the radial and tangential derivatives. Let $\partial_r V$ denote the radial derivative and $\{\partial_{s_i}V\}_{i=1}^{n-1}$ a set of mutually orthogonal tangential derivatives. Then there exists a set $M$ of measure zero such that:
\begin{align}
  &\lim_{\substack{x\in\mathbb{R}^n\setminus M \\ |x|\to\infty}} \frac{V^2(x)}{|\partial_r V(x)|} = 0, \tag{3.15}\\[4pt]
  &\exists\, C,\beta,\rho > 0:\quad C|\partial_r V(x)| > \frac{1}{|x|^\beta} \quad\text{for } |x| > \rho, \tag{3.16}\\[4pt]
  &C|\partial_r V(x)| > \Bigl(\textstyle\sum_{i=1}^{n-1}|\partial_{s_i}V(x)|^2\Bigr)^{1/2}. \tag{3.17}
\end{align}

\noindent\textbf{Interpretation.}
\begin{itemize}[nosep]
  \item (3.15) says $V$ decays faster than $|\nabla V|^{1/2}$; i.e., the potential vanishes at infinity but its gradient does not vanish ``too fast'' relative to $V$ itself.
  \item (3.16) says $|\partial_r V|$ cannot decay faster than any polynomial. This prevents $V$ from having a degenerate critical point at infinity.
  \item (3.17) says the radial component of $\nabla V$ dominates the tangential components. This is automatically satisfied for radially symmetric potentials.
\end{itemize}

\subsection{For Multi-Peak Existence and Uniqueness (Theorem~4.3)}

Let $\{x_1,\ldots,x_m\}$ be a finite subset of the critical points of $V$. The theorem requires:
\begin{enumerate}[nosep]
  \item $V$ satisfies (H1) and (H2).
  \item $V$ is \emph{twice differentiable} at each $x_j$.
  \item The Hessian $H_V(x_j)$ is \emph{nonsingular} (i.e., $x_j$ is a non-degenerate critical point).
\end{enumerate}
Under these conditions, for each choice of signs $\mu_i\in\{+1,-1\}$ and all sufficiently large $E$, equation~(1.3) has a \emph{unique} real-valued solution $U_E\in H^2$ satisfying
\[
  \bigl\|U_E - \textstyle\sum_{i=1}^m \mu_i\, u_\infty(\cdot - x_i\sqrt{E})\bigr\|_{H^2} < \varepsilon,
\]
where $u_\infty$ is the unique positive radial solution of $(-\Delta+1)u + \sigma|u|^{2p}u = 0$.

\subsection{For the Symmetry-Breaking Bifurcation (Kirr--Kevrekidis--Pelinovsky, 2011)}

The 1D result additionally assumes:
\begin{enumerate}[label=\textbf{(H\arabic*)}, start=4]
  \item $-E_0 < 0$ is the lowest eigenvalue of $L_0 = -\partial_x^2 + V(x)$.
\end{enumerate}
and, for the main bifurcation theorem:
\begin{enumerate}[label=\textbf{(H\arabic*)}, start=5]
  \item $xV'(x) \in L^\infty(\mathbb{R})$.
\end{enumerate}
The authors note that (H1) and (H5) can be relaxed to $V(x),\, xV'(x)\in L^q(\mathbb{R})$ for some $q\geq 1$.

\subsection{Summary for the Draft}

For your draft, the \texttt{\textbackslash begin\{comment\} List them \textbackslash end\{comment\}} placeholder in Section~2 should be replaced with something like:

\medskip
\noindent\emph{``Under the following conditions on the potential:}
\begin{enumerate}[nosep,label=(\roman*)]
  \item \emph{$V\in W^{1,\infty}(\mathbb{R}^n)$ with $\lim_{|x|\to\infty}V(x) = 0$ (ensures self-adjointness, compactness of $V$ relative to $-\Delta$, and Weyl's theorem for the essential spectrum);}
  \item \emph{$x\cdot\nabla V\in L^\infty(\mathbb{R}^n)$ (controls the Pohozaev identity governing the $E\to\infty$ scaling);}
  \item \emph{$V^2(x)/|\partial_r V(x)|\to 0$ as $|x|\to\infty$, with $|\partial_r V|$ dominating all tangential derivatives and decaying no faster than polynomially (conditions (3.15)--(3.17) of~\cite{KirrNatarajan2018}, preventing profile drift to infinity);}
  \item \emph{$-\Delta+V$ supports at least one bound state (eigenvalue $-E_0 < 0$),}
\end{enumerate}
\emph{the rescaled positions $E^{-1/2}z_i(E)$ remain bounded and converge to the critical points of $V$ as $R\to 0$.''}


%% ============================================================
\section{Gap Analysis: Kirr's Published Work vs.\ This Research}
%% ============================================================

The flowchart (Figure~7) delineates two regions: the framework established by Kirr--Natarajan~(2018) and the contributions of the present work. Here we analyze the precise junction between them and identify what is already covered in the literature.

\subsection{What Kirr--Natarajan (2018) Establishes}

\begin{enumerate}[nosep]
  \item \textbf{Global bifurcation framework} (Sections~2--3): Maximal $C^1$ branches of solutions, dichotomy between $E_*=\infty$ and finite bifurcation, compactness via concentration-compactness.
  \item \textbf{$E\to\infty$ limit} (Theorem~4.1): Under the rescaling $u_E(x) = E^{-1/(2p)}\psi_E(E^{-1/2}x)$, the norm scaling $\|\psi_E\|_{L^{2p+2}}^{2p+2}\sim E^{1/2+1/p}$ and convergence of $u_E$ to superpositions of the ground state $u_\infty$ of the translation-invariant equation.
  \item \textbf{Single profile at each critical point} (Theorem~4.3): For $m$ distinct non-degenerate critical points of $V$, existence and \emph{uniqueness} of a solution with one profile $u_\infty$ centered at each. Orbital stability determined by the sign of $V''(x_j)$ and the nonlinearity power $p$.
  \item \textbf{The open case} (Remark~4.1): ``We are missing the case when the superposition has multiple profiles $u_\infty$ approaching the same critical point. \ldots our current partial results show that multiple profiles at a minimum of $V$ cannot occur but saddle points and maxima of $V$ allow for any number of profiles to approach it along the directions given by the eigenvectors of the Hessian of $V$ at the saddle or maximal point corresponding to negative eigenvalues. \emph{We will analyze it in a future paper.}''
\end{enumerate}

\subsection{What the Present Work Contributes (the ``Future Paper'')}

The thesis is precisely the analysis promised in Remark~4.1 of Kirr--Natarajan. Specifically:

\begin{enumerate}[nosep]
  \item \textbf{Classification of multi-profile configurations at a single critical point}: The collinear ($M$ peaks along one eigenvector), isosceles, equilateral, and rotationally parameterized triangular configurations in the plane spanned by two eigenvectors of $H_V$.
  \item \textbf{Admissibility constraints}: The explicit $(\mu,\theta)$ regions where $\alpha_{ki}\geq 0$.
  \item \textbf{The perturbation system $F_i(\delta y, R)=0$}: Formulation and explicit solution $\delta y^0$ at $R=0$.
  \item \textbf{Kernel identification and decomposition}: Showing $\Omega = \ker(D_{\delta y}F_i)$ is nontrivial (translations), decomposing into $P_\perp F \oplus P_\parallel F$, and identifying the gauge-fixing conditions that remove the rotational degeneracy.
\end{enumerate}

\subsection{What Remains Open}

Three items are identified in the Conclusion (Section~5 of the draft):
\begin{enumerate}[nosep]
  \item Rigorous verification of rotational equivariance $F_i(R_\omega\delta y, R) = R_\omega F_i(\delta y, R)$.
  \item Invertibility of $D_{\delta y}P_{\text{rot}}^\perp F_i(\delta y^0, 0)$ on the gauge-fixed subspace.
  \item Extension to $\theta$-dependent configurations (where $\alpha_{ij}$ depend on $\theta$).
\end{enumerate}
Once completed, Theorem~4.3 of Kirr--Natarajan lifts the finite-dimensional solutions back to exact PDE ground states for all sufficiently large $E$.

\subsection{Is There Other Published Work by Kirr Covering This Gap?}

Based on the references and the explicit statement in Remark~4.1 (``we will analyze it in a future paper''), this multi-profile-at-a-single-critical-point analysis has \textbf{not} appeared in Kirr's published work as of 2018. The 2011 paper (Kirr--Kevrekidis--Pelinovsky) addresses symmetry-breaking bifurcation in 1D for symmetric (double-well) potentials, which is a different (though related) phenomenon. The survey paper Kirr~(2016) [arXiv:1605.08167] articulates the general program but does not carry out the multi-peak classification.

\textbf{Conclusion:} Your thesis fills a gap that Kirr--Natarajan explicitly identified as open. You should cite Remark~4.1 of their paper and state clearly that the present work constitutes the analysis announced therein, specialized to the case of a non-degenerate critical point with negative Hessian eigenvalues.

\end{document}
